\documentclass[10pt,a4paper]{amsart}
\usepackage[margin=19mm]{geometry}
\usepackage{amsmath,amssymb,mathtools}
\usepackage[scr=rsfs,cal=euler]{mathalfa}
\usepackage{xfrac}
\usepackage[unicode,hidelinks,bookmarks=false]{hyperref}
\newcommand{\ie}{i.e.\ }
\newcommand{\cf}{cf.\ }
\newcommand{\resp}{respectively\ }
\newcommand{\Subsection}{\subsection}
\providecommand{\nobreakdash}{\nobreak\hbox{-}}
\setlength{\emergencystretch}{2em}
\multlinegap0pt
\usepackage{bm}
\usepackage{bbm}
\usepackage{dsfont}
\usepackage{xspace}
\usepackage{cancel}
\usepackage{mathtools}
\usepackage{amsthm}
\makeatletter
\@ifundefined{lemma}{\newtheorem{lemma}{Lemma}}{}
\@ifundefined{theorem}{\newtheorem{theorem}{Theorem}}{}
\makeatother
\title[Remoteness, order, size and connectivity constraints in digr]{Remoteness, order, size and connectivity constraints in digraphs}
\author{Sufiyan Mallu}
\date{Source: arXiv:2510.08841v1 (CC BY 4.0)}
\begin{document}
\maketitle

\noindent\emph{Source and licence.} Remoteness, order, size and connectivity constraints in digraphs. Version arXiv:2510.08841v1 (CC BY 4.0), distributed under the Creative Commons Attribution 4.0 International licence, as recorded in the arXiv licence metadata for this exact version and shown on the abstract page https://arxiv.org/abs/2510.08841v1. Full rights notice: licenses/arxiv-2510.08841-CC-BY-4.0.txt.\par\smallskip

\noindent\textit{Editorial scope.} Counted unit: the Theorem whose source label is \texttt{eq:bound-on-rho-for-kappa-digraph} in the article (source0.tex lines 302--309, source label \texttt{eq:bound-on-rho-for-kappa-digraph}), delivered together with the article's own complete original proof of it (source0.tex lines 311--364, one \texttt{proof} environment of 54 lines). No mathematics is written, repaired or re-derived: statement and proof are the article's own text line for line. Required context retained uncounted: la:size-vs-remoteness-in-kappa-connected-pc-digraph (lines 250--257). Every definition, display and label of those lines is retained unchanged. Omitted from this package and not redistributed: 1--249, 258--301, 310--310, 365--702. Nothing inside a retained excerpt is omitted or altered: every retained body is delivered exactly as the article's lines are written, so no omission receipt of any kind accompanies this package. Citations: the retained excerpts carry no citation command at all, so no bibliography entry of the article is transcribed and no reference is redistributed. Reference adaptation: none was needed and none was made; every reference inside the delivered unit resolves inside it, so no unresolved reference or citation is produced. Printed slips kept literal: no printed word, symbol, hypothesis or number of the counted statement or of its proof was altered. Layout and class: the article's own class is \texttt{article} with packages including authblk, babel, amsmath, amsfonts, amsthm, amssymb, pgf, tikz, pgfplots, float, graphicx, tikz-cd; the wrapper is the lane's pinned \texttt{amsart} editorial wrapper at 10pt on A4, which restates the theorem-like environments the retained text needs and adds the article's own macro definitions that the retained text needs (none), with extra wrapper packages (bm, bbm, dsfont, xspace, cancel, mathtools). The delivered numbering is the unit's own. Assets: no retained span contains a figure environment, a graphics-inclusion command, a PDF-page inclusion, a table or a TikZ picture; no image file is copied into this package, the package declares no asset dependency and no image file is redistributed with it. Licence: version arXiv:2510.08841v1 of this article is distributed under the Creative Commons Attribution 4.0 International licence, as recorded in the arXiv licence metadata for this exact version and shown on the abstract page https://arxiv.org/abs/2510.08841v1. A full rights notice is delivered at licenses/arxiv-2510.08841-CC-BY-4.0.txt. Characters: every retained excerpt is pure ASCII, so the delivered engine, XeLaTeX with its default Latin Modern text font, reports no missing character.
\begin{lemma}\label{la:size-vs-remoteness-in-kappa-connected-pc-digraph}
	(a) Let $H$ be a $\kappa$-connected path-complete digraph, where $\kappa \in \mathbb{N}$. Then 
$\rho(H+\overrightarrow{uv}) < \rho(H)$ for any arc $ \overrightarrow{uv} \in E(\overline{H})$. \\
(b) Let $H_1, H_2$ be two distinct $\kappa$-connected path-complete strong digraphs of order $n$. Then either $m(H_1) < m(H_2)$ and $\rho(H_1) > \rho(H_2)$, or $m(H_1) > m(H_2)$ 
and $\rho(H_1) < \rho(H_2)$. \\
 (c) Given $n$, $\kappa$. Then there exists a $\kappa$-connected path-complete strong digraph of order $n$ and size $m$ if an only if $m \equiv (n^2-2n-1) \pmod{\kappa} $ and $ \frac{n}{2}(3\kappa+n)-n-\kappa^2-b(\kappa-b) \leq m \leq n^2-2n-1$, where $b \in \{1,2,\ldots, \kappa\}$ with
$b \equiv n-1 \pmod{\kappa}$.
\end{lemma}

	\begin{theorem}
(a) Let $D$ be a $\kappa$-connected strong digraph of order $n$ and size $m$ with $ m \leq n^2-2n-1$. Then 
	\begin{equation} \label{eq:bound-on-rho-for-kappa-digraph}
\rho(D) \leq \rho(DPK_{n,m,\kappa}).
\end{equation}
(b) Assume that $m \equiv (n^2-2n-1) \pmod{\kappa} $ and $ \frac{n}{2}(3\kappa+n)-n-\kappa^2-b(\kappa-b) \leq m \leq n^2-2n-1$, where $b$ is the integer in $\{1,2,\ldots, \kappa\}$ with
$b \equiv n-1 \pmod{\kappa}$. Then equality in (a) holds only if $D=DPK_{n,m,\kappa}$.
	\end{theorem}

	\begin{proof}
We first prove that there exists a $\kappa$-connected path-complete $(n,m)$- strong digraph $D'$ with 
\begin{equation} \label{eq:exists-pc-kappa-digraph}
\rho(D) \leq \rho(D'). 
\end{equation}
We may assume that $D$ has maximum remoteness among all $\kappa$-connected $(n,m)$- strong digraphs, and that among all such strong digraphs with maximum remoteness, $D$ is one with the maximum size. Furthermore, let $v \in V(D)$ with $\overline{\sigma}(v,D) = \rho(D)$, $d = {\rm ecc}_D(v)$, $N_{i} = \{z \in V(D) \mid d_{D}(v,z) = i \}$, and $|N_{i}| = n_i$ for $i \in \{0, 1, \ldots, d\}$. Clearly, $n_0 = 1$,  and $\sum_{i=0}^{d} n_i = n$. \\

\emph{{\bf Claim 1:}} $D=K_{n_0} \; \overleftarrow{+} \; \overleftrightarrow{K_{n_1}} \; \overleftarrow{+} \ldots \; \overleftarrow{+} \;  \overleftrightarrow{K_{n_{d-1}}} \; \overleftarrow{+} \; \overleftrightarrow{K_{n_{d}}}$. \\
 
Recall that \( D \) has the maximum size among all strong digraphs of size at least \( m \) for which \( \sigma(v, D) \) is maximised. Consequently, each \( N_i \) induces a complete digraph \( \overleftrightarrow{K_{n_i}} \) in the digraph \( D \), otherwise we could add arcs between vertices of $N_i$ without changing the remoteness of $D$. Additionally, \( D \) has arcs from every vertex in \( N_i \) to every vertex in \( N_{i+1} \), and arcs from every vertex in \( N_{i+1} \) to every vertex in \( N_i \) for \( i = 0, 1, 2, \ldots, d-1 \). Furthermore, every vertex in \( N_j \) is adjacent to every vertex in \( N_i \) for all \( j > i \) where \( i, j \in \{0, 1, \ldots, d\} \). Hence, Claim 1 follows. \\

Note that $n_i \geq \kappa$ holds for $i=1,2,\ldots,d-1$, but not for $0$.\\

\emph{{\bf Claim 2:}} For all $ i \in \{1,2, \ldots, d-3\} $, we have $ n_i = \kappa $. \\

Suppose to the contrary that there exists $ j \in \{1,\ldots,d-3\} $ with $ n_{j}>\kappa $. Let $j$ be the smallest such value. Then $n_{i} = \kappa$ for all $ i \in \{1,\ldots,j-1\} $. Now consider the strong digraph $D^*$ that is obtained from $D$ by moving a vertex from $N_j$ to $N_{j+1}$, i.e., $D^*=K_{n'_0} \; \overleftarrow{+} \; \overleftrightarrow{K_{n'_1}} \; \overleftarrow{+} \ldots \; \overleftarrow{+} \; \overleftrightarrow{K_{n'_d}}$, where
$n'_i=n_i$ for $i \in \{0,1,\ldots,d\} - \{j,j+1\}$,  
$n'_j=n_j-1$ and $n'_{j+1}= n_{j+1}+1$. Then $m(D^*) = m(D) + 2(n_{j+2} - n_{j-1}) \geq m(D)$ since $n_{j+2} \geq \kappa=n_{j-1}$ and $D^{*}$ is $(n,m)$-strong digraph with $\sigma(v,D^{*}) > \sigma (v,D)$, and thus $\rho(D^{*}) > \rho(D)$. This contradiction to the maximality of $\rho(D)$ proves Claim 2.  \\

\emph{{\bf Claim 3:}} $ n_{d-2} = \kappa $. \\

Since $D$ is $k$-connected, we have that $n_{d-2} \geq \kappa$. Now, suppose to the contrary that \( n_{d-2} \neq \kappa \). Then \( n_{d-2} > \kappa \). Consider the strong digraph \( D^* \), which is obtained from \( D \) by moving one vertex from \( N_{d-2} \) and \( N_d \) to \( N_{d-1} \), i.e., $D^*=K_{n'_0} \; \overleftarrow{+} \; \overleftrightarrow{K_{n'_1}} \; \overleftarrow{+} \ldots \; \overleftarrow{+} \; \overleftrightarrow{K_{n'_d}}$, where
$n'_i=n_i$ for $i \in \{0,1,\ldots,d-3\} - \{d-2,d-1,d\}$,  
$n'_{d-2}=n_{d-2}-1$, $n'_{d}=n_{d}-1$ and $n'_{d-1}= n_{d-1}+2$. It is easy to verify that $m(D^*) = m(D) + 2n_d - n_{d-3} + n_{d-2} - 1$. Since \( n_{d-3} = \kappa \) and \( n_{d-2} \geq \kappa+1 \), we have $
m(D^*) \geq m(D) + 2n_d-1 > m(D).
$ Moreover, \( \rho(D^*) = \rho(D) \). This contradicts our choice of $D$ as a digraph with the maximum size among those of maximum remoteness. Thus, Claim 3 follows.  \\



\noindent It follows from Claims 1 to 3 that $D$ is a $\kappa$-connected path-complete strong digraph. Letting $D'=D$ proves (\ref{eq:exists-pc-kappa-digraph}).  \\

By (\ref{eq:exists-pc-kappa-digraph}), there exists a $\kappa$-connected path-complete strong digraph $D'$ of order $n$ and size at least $m$ with $\rho(D) \leq \rho(D')$. By the definition of $DPK_{n,m,\kappa}$, we have $m(D') \geq m(DPK_{n,m,\kappa})$.  By Lemma \ref{la:size-vs-remoteness-in-kappa-connected-pc-digraph}(a), it follows that $\rho(D) \leq \rho(DPK_{n,m,\kappa}) $. Hence 
\[\rho(D) \leq \rho(D') \leq \rho(DPK_{n,m,\kappa}),\]
which proves (a). \\

(b) Now assume that equality holds in \eqref{eq:bound-on-rho-for-kappa-digraph}, i.e., that
$\rho(D) = \rho(DPK_{n,m,\kappa})$, and furthermore that $m \equiv (n^2-2n-1) \pmod{\kappa} $ and $ \frac{n}{2}(3\kappa+n)-n-\kappa^2-b(\kappa-b) \leq m \leq n^2-2n-1$, where $b$ is as defined above. 
It follows from Lemma \ref{la:size-vs-remoteness-in-kappa-connected-pc-digraph}(c) that the 
graph $DPK_{n,m,\kappa}$ has exactly $m$ arcs. \\

It follows from part (b) that $D$ has maximum remoteness among all $\kappa$-connected 
$(n,m)$-strong digraphs. We claim that $D$ has maximum size among all such strong digraphs 
maximising the remoteness. Suppose not. Then there exists a $\kappa$-connected 
$(n,m+1)$-strong digraph $D''$ with $\rho(D) = \rho(D'')$. Applying \eqref{eq:bound-on-rho-for-kappa-digraph} to $D''$ we get that
\[ \rho(D) = \rho(D'') \leq \rho(DPK_{n,m+1,\kappa}) < \rho(DPK_{n,m,\kappa}), \]
where the last inequality follows from 
Lemma \ref{la:size-vs-remoteness-in-kappa-connected-pc-digraph}(b) and the fact that 
$m(DPK_{n,m,\kappa})= m < m(DPK_{n,m+1,\kappa})$. Hence $D$ has maximum size among all $\kappa$-connected strong digraphs of order $n$ maximising the remoteness. \\

The proof of (a) shows that, if $D$ has maximum size among all $\kappa$-connected
path-complete $(n,m)$-strong digraphs, then $D$ is a path-complete strong digraph. Hence $D=DPK_{n,m',\kappa}$ for some $m'$ with $m' \geq m$. Since by Lemma \ref{la:size-vs-remoteness-in-kappa-connected-pc-digraph} we have 
$\rho(DPK_{n,m',\kappa}) < \rho(DPK_{n,m,\kappa})$ if $m'>m$, we have $m'=m$, 
and thus $D=DPK_{n,m,\kappa}$, as desired. 
\end{proof}

\end{document}
